\documentclass[12pt]{article}
\usepackage[utf8]{inputenc}
\usepackage[T1]{fontenc}
\usepackage{hyperref}
\title{Utu Language Resources for K-12 Curriculum Development}
\author{Compiled by Claude Sonnet 5 (claude-sonnet-5) for Andr\'as Kornai}
\date{}
\begin{document}
\maketitle

\section{Introduction}
Utu (ISO 639-3: utu) is a Trans-New Guinea (Madang) language of Papua New Guinea, spoken about 30 km from Madang town. Do not confuse it with Uya (Usu), a separate Rai Coast language. Very little was found: general profiles only, with no teaching or reference materials.

\section{Primary Educational Resources}
None found.

\section{Linguistic and Reference Materials}
\begin{itemize}
\item Wikipedia overview of the Utu language. \url{https://en.wikipedia.org/wiki/Utu_language}
\item Endangered Languages Project profile. \url{https://endangeredlanguages.com/lang/3409}
\item ASJP entry for Utu. It cites Z'graggen's 1980 comparative word list of the Northern Adelbert Range languages (Pacific Linguistics D-32), which I did not locate online. \url{https://asjp.clld.org/languages/UTU}
\end{itemize}

\section{Religious and Community Resources}
Joshua Project reports that the Utu and their neighbouring language groups have no scripture in their own languages. \url{https://joshuaproject.net/people_groups/15751}

\section{Digital and Audio Resources}
None found.

\section{Curriculum Development Context}
Wikipedia cites 1,010 speakers (2000 census), and Joshua Project says all generations use the language face to face. No orthography, literacy materials or school programme were found. Any curriculum would require new documentation work with the community.
\end{document}